\chapter{Ice-Breaker Không Lố}
\chapterintro{Khởi động vừa đủ, không biến tiết học thành sân khấu}{Màn mở vui chỉ có giá trị nếu kéo được lớp vào bài, không phải kéo bài ra khỏi lớp.}

\section{Giơ tay theo cảm giác}
\begin{tinhhuong}{Tình huống}
Muốn khởi động nhanh mà không cần phát biểu dài dòng.
\end{tinhhuong}
\begin{goiydungngay}
Đặt một câu đơn giản: \textit{``Ai thấy mình đang tỉnh, giơ tay; ai thấy cần cứu, giơ hai tay.''}
\begin{itemize}
\item Cả lớp tham gia trong 5 giây.
\item Không khí bật lên nhanh.
\item Sau đó nối ngay vào bài.
\end{itemize}
\end{goiydungngay}
\begin{cauchot}
Ice-breaker tốt là thứ mở lớp, không nuốt lớp.
\end{cauchot}
\ngancach

\section{Một lựa chọn vui}
\begin{tinhhuong}{Tình huống}
Lớp cần bật nói nhưng không muốn dùng câu hỏi quá dài.
\end{tinhhuong}
\begin{goiydungngay}
Cho chọn nhanh giữa hai đáp án vui: \textit{``Bài hôm nay hợp kiểu cà phê hay trà sữa?''}
\begin{itemize}
\item Chỉ cần chọn và nói 1 lý do.
\item Không sợ sai.
\item Tạo tiếng cười nhẹ rất nhanh.
\end{itemize}
\end{goiydungngay}
\begin{cauchot}
Câu hỏi nhẹ là chiếc cầu tốt nhất vào bài nặng.
\end{cauchot}
\ngancach

\section{Nói nối tiếp 3 từ}
\begin{tinhhuong}{Tình huống}
Muốn cả lớp nhập cuộc cùng lúc mà không cần nhiều chuẩn bị.
\end{tinhhuong}
\begin{goiydungngay}
Yêu cầu 3 em nối tiếp nhau, mỗi em chỉ được nói 1 từ để tạo thành một câu vui mở bài.
\begin{itemize}
\item Rất nhanh.
\item Cả lớp dễ cười.
\item Hợp đầu giờ hoặc sau giờ giải lao.
\end{itemize}
\end{goiydungngay}
\begin{cauchot}
Mở lớp nhanh không cần đạo cụ, chỉ cần đúng nhịp.
\end{cauchot}
\ngancach

\section{Bốc thăm bằng cử chỉ}
\begin{tinhhuong}{Tình huống}
Muốn mở lớp vui nhưng không muốn gọi tên học sinh quá ngẫu nhiên làm các em căng.
\end{tinhhuong}
\begin{goiydungngay}
Quy ước nhanh: ai đang cầm bút xanh nói trước, ai đang đeo đồng hồ nói sau, ai ngồi cạnh cửa nói tiếp.
\begin{itemize}
\item Cách chọn người có chút bất ngờ nhưng không đáng sợ.
\item Cả lớp bật chú ý rất nhanh.
\item Dùng xong là vào bài ngay.
\end{itemize}
\end{goiydungngay}
\begin{cauchot}
Một chút ngẫu nhiên nhẹ nhàng giúp lớp tỉnh mà không lố.
\end{cauchot}
\ngancach

\section{Chọn biểu tượng để mở bài}
\begin{tinhhuong}{Tình huống}
Lớp cần một câu mở vui nhưng phải đủ gọn để không lấn thời gian dạy.
\end{tinhhuong}
\begin{goiydungngay}
Vẽ nhanh 3 biểu tượng lên bảng như mặt cười, đám mây, tia chớp rồi hỏi lớp bài hôm nay giống biểu tượng nào nhất.
\begin{itemize}
\item Học sinh chọn rất nhanh.
\item Thầy cô dễ nối biểu tượng sang tinh thần bài.
\item Không cần chuẩn bị đạo cụ hay slide.
\end{itemize}
\end{goiydungngay}
\begin{cauchot}
Ice-breaker càng đơn giản, cơ hội nối mượt vào bài càng cao.
\end{cauchot}
